\documentclass[10pt,a4paper]{amsart}
\usepackage[margin=19mm]{geometry}
\usepackage{amsmath,amssymb,mathtools}
\usepackage[scr=rsfs,cal=euler]{mathalfa}
\usepackage{xfrac}
\usepackage[unicode,hidelinks,bookmarks=false]{hyperref}
\newcommand{\ie}{i.e.\ }
\newcommand{\cf}{cf.\ }
\newcommand{\resp}{respectively\ }
\newcommand{\Subsection}{\subsection}
\providecommand{\nobreakdash}{\nobreak\hbox{-}}
\setlength{\emergencystretch}{2em}
\multlinegap0pt
\usepackage{amssymb}
\usepackage[shortlabels]{enumitem}
\newtheorem{proposition}{Proposition}
\newcommand{\Q}{\mathbb{Q}}
\newcommand{\Z}{\mathbb{Z}}
\title{Automorphic Cohomology and the Limits of Algebraic Cycles}
\author{Amir Mostaed}
\date{Source: arXiv:2602.06865v1 (CC BY 4.0)}
\begin{document}
\maketitle

\noindent\emph{Source and licence.} Automorphic Cohomology and the Limits of Algebraic Cycles. Version arXiv:2602.06865v1 (CC BY 4.0), distributed by arXiv under the Creative Commons Attribution 4.0 International licence, http://creativecommons.org/licenses/by/4.0/, as recorded in the arXiv licence metadata for this exact version and shown on the abstract page https://arxiv.org/abs/2602.06865v1. Full licence notice: \texttt{licenses/arxiv-2602.06865-CC-BY-4.0.txt}.\par\smallskip

\noindent\textit{Editorial scope.} Counted unit: the article's proposition proposition (source0.tex lines 853--855). The article's complete original proof of it is delivered with it (source0.tex lines 857--902, one \texttt{proof} environment of 46 lines). No mathematics is written, repaired or re-derived: the statement and the proof are the article's own lines, in the article's own notation, and no retained line is substituted or re-flowed.  No further source-local context is needed: every cross-reference of every retained line resolves inside the two delivered spans.  Nothing was omitted from any retained span, so the package carries no omission record.  No retained line contains a citation, so the wrapper carries no bibliography and no bibliography entry.  No retained line contains a figure environment, an image-inclusion command or drawing code, so no image is delivered and the package declares no asset dependency.  Editorial formatting only, in addition: an AMS \texttt{amsart} preamble that restates the article's own declarations and macros used by the retained lines, namely the environments \texttt{proposition} on separate counters, together with the standard packages those lines need.  The retained environments print in this document as Proposition 1 in order of appearance, whereas the article prints the same items with the numbering of its own full document.  The complete native source of the article as distributed in the arXiv e-print for this exact version is bundled unchanged as \texttt{sources/194128/source0.tex}.
\begin{proposition}[Hecke field is $\Q$]
For $f \in S_2(\Gamma_0(11))$, since $\dim S_2(\Gamma_0(11)) = 1$, the form $f$ is unique and rational ($a_n \in \Z$ for all $n$). Therefore the Hecke field is $E = \Q(a_2, a_3, \ldots) = \Q$ with $[E : \Q] = 1$.
\end{proposition}

\begin{proof}
We prove this in several steps, combining the theory of modular forms, elliptic curves, and representation theory. The genus of the modular curve $X_0(11)$ is computed via the Riemann--Hurwitz formula. For $X_0(N)$, the genus is $g(X_0(N)) = 1 + \frac{\mu}{12} - \frac{\nu_2}{4} - \frac{\nu_3}{3} - \frac{\nu_\infty}{2},$ where 
$\mu = [\mathrm{SL}_2(\Z) : \Gamma_0(N)]$ is the index, $\nu_2$ = number of elliptic points of order 2, $\nu_3$ = number of elliptic points of order 3, $\nu_\infty$ = number of cusps. For $N = 11$ (prime),
\begin{align*}
\mu &= N \prod_{p \mid N} \left(1 + \frac{1}{p}\right) = 11 \cdot \left(1 + \frac{1}{11}\right) = 12, \\
\nu_2 &= 0 \quad \text{(no elliptic points of order 2 for } N = 11), \\
\nu_3 &= 0 \quad \text{(no elliptic points of order 3 for } N = 11), \\
\nu_\infty &= 2 \quad \text{(two cusps: } 0 \text{ and } \infty).
\end{align*}

Therefore $g(X_0(11)) = 1 + \frac{12}{12} - 0 - 0 - \frac{2}{2} = 1 + 1 - 1 = 1$. By the Riemann--Roch theorem for modular forms, the dimension of $S_2(\Gamma_0(N))$ equals the genus $\dim S_2(\Gamma_0(N)) = g(X_0(N))$. Therefore $\dim S_2(\Gamma_0(11)) = 1$. Since $\dim S_2(\Gamma_0(11)) = 1$, there is a unique normalized eigenform (up to scaling)
$f(z) = \sum_{n=1}^{\infty} a_n q^n \quad \text{with } a_1 = 1$. This $f$ is automatically a newform (not coming from lower level) because 11 is prime. By the modularity theorem (Wiles--Taylor--Diamond), every elliptic curve over $\Q$ corresponds to a weight-2 newform. Conversely, by work of Shimura and Taniyama, a weight-2 newform of level $N$ with rational Fourier coefficients corresponds to an elliptic curve of conductor $N$. The elliptic curve of conductor 11 is unique (up to isogeny) $E_{11}: y^2 + y = x^3 - x^2.$ This is the minimal Weierstrass equation (Cremona label 11a3). For an elliptic curve $E/\Q$ given by a Weierstrass equation with coefficients in $\Q$, the number of points over $\mathbb{F}_p$ is $\#E(\mathbb{F}_p) = p + 1 - a_p$ where $a_p \in \Z$ is the $p$-th Fourier coefficient of the associated modular form. Since $E_{11}$ is defined over $\Q$ (with equation in $\Q$), we have $\#E_{11}(\mathbb{F}_p) \in \Z$ \text{for all primes} $p \neq 11$.  This immediately gives $a_p = p + 1 - \#E_{11}(\mathbb{F}_p) \in \Z$.

For $E_{11}: y^2 + y = x^3 - x^2$, the computation is as follows.

\textit{At $p = 2$:}
Over $\mathbb{F}_2$, the curve equation becomes $y^2 + y = x^3 + x^2$. Counting points:
$$\#E_{11}(\mathbb{F}_2) = 5.$$
Therefore
$$a_2 = 2 + 1 - 5 = -2 \in \Z.$$

\textit{At $p = 3$:}
Over $\mathbb{F}_3$, counting points gives
$$\#E_{11}(\mathbb{F}_3) = 5.$$
Therefore
$$a_3 = 3 + 1 - 5 = -1 \in \Z.$$

\textit{At $p = 5$:}
$$\#E_{11}(\mathbb{F}_5) = 5.$$
Therefore
$$a_5 = 5 + 1 - 5 = 1 \in \Z.$$

\textit{At $p = 7$:}
$$\#E_{11}(\mathbb{F}_7) = 10.$$
Therefore
$$a_7 = 7 + 1 - 10 = -2 \in \Z. $$

At the prime of bad reduction $p = 11$, the curve has multiplicative reduction. The reduction type is split multiplicative (Tate curve), giving $a_{11} = \pm 1$.
The sign is determined by the reduction being split or non-split. For $E_{11}$, the reduction is split, giving
$a_{11} = 1 \in \Z$.
The Fourier coefficients satisfy the recurrence relations from the Hecke operators, $a_{mn} = a_m a_n$ \text{for} $\gcd(m,n) = 1, a_{p^{k+1}} = a_p a_{p^k} - p a_{p^{k-1}}$  \text{for primes} $p \neq 11$, and 
$a_{11^k} = a_{11}^k = 1$ \text{(Steinberg)}. Starting from $a_2, a_3, a_5, a_7, \ldots \in \Z$, these recurrence relations imply
$a_n \in \Z \quad \text{for all } n \geq 1$. The Hecke field is defined as
$E = \Q(a_2, a_3, a_5, a_7, \ldots) = \Q(a_n : n \geq 1)$. Since $a_n \in \Z$ for all $n$, we have
$E = \Q$. Therefore $[E : \Q] = 1$.
\end{proof}

\end{document}
